\section{Join the coefficient path to the pixel path}
A local reconstruction must follow the representation the decoder receives.
The forward transform's coefficients describe the source block, but quantization
has already replaced them with selected levels. Reconstruct those levels before
calling the inverse. For an INTER block, add its inverse residual to the same
prediction used to form the original residual. For an INTRA block, the inverse
already describes absolute samples; its prediction is zero.

Write the current sample as $S_i$, prediction as $P_i$ and inverse output as $D_i$.
The transform input is $S_i-P_i$ for INTER and $S_i$ for INTRA. The reconstructed
pixel is
\[
 X_i=\begin{cases}
 \operatorname{clip}_{[0,255]}(D_i)&\text{INTRA},\\
 \operatorname{clip}_{[0,255]}(P_i+D_i)&\text{INTER}.
 \end{cases}
\]
The addition precedes final sample clipping. If prediction is 250 and the signed
inverse residual is 20, their sum is 270 and the stored sample is 255. A byte
addition would wrap to fourteen. If prediction is five and the residual is
minus ten, clipping the residual to zero first would incorrectly keep five;
the complete signed sum is minus five and the correct stored sample is zero.

The inverse's signed output lies in $[-256,255]$, and prediction lies in
$[0,255]$. Their sum therefore lies in $[-256,510]$, which fits a signed halfword.
This bound permits ordinary halfword vector addition followed by signed minimum
and maximum operations. It does not permit byte wrapping. Store pixels as bytes
only after clipping the signed sum.

\section{Follow a constant block through every representation}
Use current samples 130, prediction samples 128 and $Q=8$. The constant residual
is two. Its forward DC is sixteen and all AC terms are zero. The exercise's
ordinary quantizer chooses level one, which reconstructs as coefficient 23.
An accurate inverse produces three at each position, so the final block contains
131. Local reconstruction has changed the source by one even though prediction,
transform and arithmetic were implemented correctly.

\begin{center}
\tagpdfsetup{table/header-rows={1}}
\begin{tabular}{lll}
\textbf{Boundary} & \textbf{First value} & \textbf{Representation}\\
Current input & 130 & Unsigned sample byte\\
Prediction & 128 & Unsigned sample byte\\
Transform input & 2 & Signed residual halfword\\
Forward DC & 16 & Transform coefficient\\
Selected DC level & 1 & Ordinary INTER level\\
Reconstructed DC & 23 & Inverse-transform coefficient\\
Inverse residual & 3 & Signed spatial halfword\\
Reconstructed sample & 131 & Clipped unsigned pixel
\end{tabular}
\end{center}
This named trace is an observation of the tested transform implementations.
Other accurate transforms need not produce identical results at every rounding
boundary. Behavioral checks must allow their documented accuracy while ensuring
that each implementation reconstructs its own selected levels consistently.

For the next block, let current samples be 132 and use the preceding reconstruction
131 as prediction. Residual one has DC eight, which falls into the dead zone at
$Q=8$. All levels become zero, the inverse residual is zero, and the new
reconstructed block remains 131. An encoder can retain a visible error while
still keeping its local decoder synchronized.

Predicting that second block from the old raw 130 gives a different path:
residual two, DC sixteen, level one and inverse residual three. A local encoder
that adds three to raw 130 obtains 133, while a decoder holding reconstructed
131 adds three and obtains 134. The loss itself is expected; this disagreement
between references is an implementation error.

\section{Make a complete block result inspectable}
Implement \texttt{h263\_encode\_block8}. Its name describes the transform and
local reconstruction stage; it does not emit a bitstream. Pass an eight-byte-aligned
64-byte request in \texttt{\$a0} and a sixteen-byte-aligned 448-byte result in
\texttt{\$a1}. Request fields are unsigned native-endian doublewords:
\begin{center}
\tagpdfsetup{table/header-rows={1}}
\begin{tabular}{rll}
\textbf{Offset} & \textbf{Field} & \textbf{Meaning}\\
0, 8 & Current pointer, stride & Eight active bytes per row\\
16, 24 & Prediction pointer, stride & INTER only; both zero for INTRA\\
32 & Quantizer parameter & 1 through 31\\
40 & Block type & 0 INTER; 1 INTRA\\
48 & Arithmetic mode & 0 scalar; 1 MSA\\
56 & Reserved & Zero
\end{tabular}
\end{center}
There are eight rows. Active strides range from eight through 4096 bytes,
independently; pointers need no alignment beyond byte addressability. The two
read-only inputs may alias. All mapped extents must be non-wrapping, and output
is disjoint from inputs and request. INTRA prediction fields are both zero and
no prediction samples are read. Supply a sixteen-byte-aligned stack with private
space. Preserve saved general registers, the incoming stack pointer, input bytes
and request bytes. Vector registers are caller-saved.

Validate every record, extent, stride and full-width field before input samples.
Invalid metadata returns status one and count zero without sample reads or result
writes. A failed transform or quantizer child returns status two and count zero;
input reads and private work already performed are permitted, but the entire
caller result remains unchanged. Preserve the incoming return address across
nested calls.

Successful output contains four consecutive arrays:
\begin{center}
\tagpdfsetup{table/header-rows={1}}
\begin{tabular}{rrl}
\textbf{Offset} & \textbf{Bytes} & \textbf{Meaning}\\
0 & 128 & Selected levels, including separate INTRA DC code\\
128 & 128 & Reconstructed coefficients from those levels\\
256 & 128 & Signed inverse spatial output\\
384 & 64 & Reconstructed pixel bytes
\end{tabular}
\end{center}
Halfwords are native-endian and in natural row-major order. Return status zero
and the nonzero ordinary-level count from quantization. INTRA DC is excluded
from that count. Stage all arrays privately and publish them together only after
every dependency succeeds. The arrays expose each boundary during construction;
a later production path may consume staged coefficients and pixels without
retaining this complete teaching record.

\section{Retain accuracy without requiring private transform bits}
Carry your completed transform tables and \texttt{h263\_transform8} and
\texttt{h263\_quantize8} entries into
\path{construction/h263-reconstruction-construction/main.asm}, replacing the dependency
stubs. Retain their helpers and required data. Check the prior standalone contracts
and the inverse accuracy gate before using those routines here.
For the provided working formats, the composition helper carries your own entries
and their data into a new file. Keep helpers after their entry labels. It preserves
the block caller and entry before the dependency marker and refuses an existing
output. Custom source layouts can be joined manually; the final interface does
not require the helper.
\begin{verbatim}
python3 construction/h263-reconstruction-construction/compose.py \
    --source my-block.asm --transform my-transform.asm \
    --quantizer my-quantizer.asm --output joined-block.asm
python3 construction/h263-reconstruction-construction/prepare.py \
    --source joined-block.asm --current 130 --prediction 128 \
    --quant 8 --intra 0 --mode 1 --output block-trace.asm
\end{verbatim}
Preparation changes the supplied constant samples and request controls, retaining
all instructions. To observe INTRA reconstruction of 128, select current 128
and INTRA one. Both tools require new output paths.

The public block checker computes an independent floating forward reference.
For each coefficient, it allows the selected level to come from a coefficient
within the earlier one-unit accuracy allowance. It then checks reconstructed
coefficients exactly against the actual returned levels. It calculates an
independent inverse from that reconstructed block and permits one-unit spatial
accuracy, while requiring final pixels to equal the clipped sum of the actual
returned inverse and prediction. Thus alternative transform rounding does not
become a false private-bit requirement, and a self-consistent but wrong scale
cannot evade the independent transform checks.

This integration check supplements the earlier function contracts and full
inverse accuracy gate. It does not replace them or establish full-stream quality.
The complete encoder will later need an independent decoder and inter-picture
checks using reconstructed references. A single correct block is not proof that
an entire reconstructed picture is committed at the right time.
